% !TeX program = xelatex
\documentclass{bjutthesis}
\usepackage{placeins}
\usepackage{listings}
\usepackage[version=4]{mhchem}
\usepackage{tikz}
\usetikzlibrary{arrows.meta,positioning}
\graphicspath{{figure/}{./}}

% 以下为代码示例的全局样式，不影响普通正文。
\definecolor{codebg}{RGB}{247,248,250}
\definecolor{codecomment}{RGB}{55,125,75}
\lstset{
  basicstyle=\small\ttfamily,
  keywordstyle=\color{black},
  commentstyle=\color{codecomment},
  backgroundcolor=\color{codebg},
  frame=single,
  framerule=0.4pt,
  rulecolor=\color{black!35},
  numbers=left,
  numberstyle=\scriptsize,
  numbersep=8pt,
  breaklines=true,
  showstringspaces=false,
  columns=fullflexible
}

\makeatletter
% 浮动体与公式间距在文档开始时统一设置，避免在各章手工插入空白。
\AtBeginDocument{%
  \raggedbottom
  \setlength{\textfloatsep}{4pt plus 2pt minus 2pt}
  \setlength{\intextsep}{4pt plus 2pt minus 2pt}
  \setlength{\floatsep}{4pt plus 2pt minus 2pt}
  \setlength{\abovecaptionskip}{4pt}
  \setlength{\belowcaptionskip}{0pt}
  \setlength{\@fptop}{0pt}
  \setlength{\@fpsep}{8pt}
  \setlength{\@fpbot}{0pt}
  \setlength{\jot}{10pt}
}
\makeatother

\begin{document}
% 封面信息统一来自 ezinfo.cfg：外封、内封、声明页。
\makecover
\maketitle
\makestate

\frontmatter
% 前置部分使用罗马数字页码，并生成中英文目录。
% !TeX root = ../thesis.tex
\begin{cabstract}
本文档用于验证和展示北京工业大学硕士学位论文模板的排版能力。正文集中呈现标题层级、正文行高、数学与物理符号、量和单位、公式、三线表、矢量图、算法、程序代码以及多种参考文献类型。模板正文采用小四号字，基线间距设为 20 bp；该参数用于近似学校 Word 规范中的 1.3 倍行距，同时提供整页连续文本，便于使用者与 Word 输出进行视觉比对。示例还说明可调整参数的位置和编译检查方法，使模板既能作为最小可运行工程，也能作为排版回归测试文档。

\ckeywords{学位论文；LaTeX 模板；排版规范；行距；科技写作}
\end{cabstract}

% !TeX root = ../thesis.tex
\begin{eabstract}
This document is a visual and functional test of the Beijing University of Technology master's thesis template. It demonstrates heading levels, body-text leading, mathematical and physical symbols, quantities and units, equations, three-line tables, vector graphics, algorithms, source code, and several bibliography entry types. The normal text uses a 12 bp font with a 20 bp baseline skip to approximate the visual result of the 1.3-line-spacing requirement specified for Microsoft Word. A full-page text specimen and reproducible build instructions are included so that users can compare output and tune the template when institutional requirements change.

\ekeywords{thesis template; LaTeX; typography; line spacing; technical writing}
\end{eabstract}

\makebitoc

\mainmatter
% 正文从阿拉伯数字第 1 页重新计数。
% !TeX root = ../thesis.tex
\chapter{文字层级与版面间距}
\label{chap:introduction}

\section{模板的验证目标}
学位论文模板不仅要给出命令，还应让使用者直观看到命令的最终效果。本章把标题、正文、段落和行距集中在连续页面中，以便检查字体是否正确嵌入、标题层级是否清晰、段前段后距离是否稳定，以及同一字号在不同软件中的视觉密度是否接近。

\section{标题层级}
章标题采用三号黑体居中；一级节标题采用四号黑体；二级节标题采用小四号黑体；更低层级分别使用小四号楷体。标题前后的距离由类文件统一控制，不建议在正文中用空行或手工的 \verb|\vspace| 调整。

\subsection{二级标题示例}
二级标题用于组织同一节内相互关联的论述。标题后第一段与后续段落采用一致的首行缩进，避免依赖编辑器自动插入空白字符。

\subsubsection{三级标题示例}
三级标题适合标识局部方法或分析步骤。若标题层级过深，通常意味着章节结构可以重新组织。

\subsubsubsection{四级标题示例}
四级标题为行内标题，适合简短条目。正式论文应尽量保持各章层级深度一致。

\section{正文行高与 Word 多倍行距}
字号、行高和行距是三个不同概念。字号描述字形的标称尺寸；行高或基线间距描述相邻两行基线之间的距离；行距则是字形上下边界之间剩余的空白。Word 的“单倍行距”通常已经大于一倍字号，“多倍行距 1.3”是在该内部基准上继续放大，而不是简单计算为字号的 1.3 倍。LaTeX 的 \verb|\linespread| 同样会乘在字体预设的基线距离上，因此不能把两个软件中的倍率直接视为同一个物理量。关于这些概念和换算关系的进一步说明见文献~\cite{wdlobin2026typesetting}。

本模板没有继续叠加全局 \verb|\linespread|，而是在 \verb|\normalsize| 中把小四号正文设置为 12 bp 字号和 20 bp 基线间距，即 \verb|\@setfontsize\normalsize{12bp}{20bp}|。比值约为 1.67。这个结果用于逼近学校规范中 Word 多倍行距 1.3 的视觉密度；不同 Word 版本、字体度量和打印驱动仍可能产生细微差异，因此应使用同一段文字进行实页比较。

\subsection{整页连续文本检验}
下面的连续文字用于观察一页 A4 纸范围内的基线稳定性。排版比较时，应关闭 PDF 阅读器的“适合宽度”缩放，分别以 100\% 比例查看 PDF 与 Word 导出的 PDF。比较对象应使用同一纸张大小、页边距、字体、字号和段落缩进。不要只比较屏幕截图中两行文字的距离，因为截图缩放和抗锯齿会改变观感。

科技论文的正文通常由定义、推理、证据和结论构成。清晰的段落应围绕一个中心问题展开：开头说明本段讨论的对象，中间给出必要的数据或推导，末尾说明这些信息如何支持本节论点。排版系统负责建立一致的视觉节奏，但不能替代内容组织。稳定的行高使读者在长行之间移动视线时不易串行，适当的段落长度则帮助读者识别论证边界。

比较行距时还要关注中英文混排。中文字符大多接近方形字面，英文小写字母具有升部和降部，公式中的积分号、上下标及分式还可能超过正文常规字面范围。固定基线距离必须为这些元素保留足够空间。如果行距过小，多行公式、带上下标的变量或加粗标题会显得拥挤；如果行距过大，段落内部的联系会减弱，页面也会出现不必要的疏松感。

段前段后距应与行高分开判断。本模板正文段落不额外增加段间距，而用首行缩进表示新段落；标题则通过类文件中的 \verb|\titlespacing| 和 \verb|\ctexset| 设置上下距离。图表与正文之间的距离由 \verb|\textfloatsep|、\verb|\intextsep| 和 \verb|\floatsep| 控制。多行公式内部的额外距离由 \verb|\jot| 控制。这些参数职责不同，混用会导致短页面看似正常、长页面却出现不一致的空白。

实际校准建议以学校提供的 Word 规范页作为参照。先在 Word 中使用指定字体、小四号和多倍行距 1.3，导出 PDF；再把同一段纯文本放入本节，分别测量二十行或三十行文字从第一条基线到最后一条基线的总高度。若累计差异明显，只调整类文件中 \verb|\normalsize| 的第二个参数，例如在 19.5 bp 至 20.5 bp 范围内试验，而不要同时修改字号和段落距离。

当正文中出现公式、图表或列表时，页面节奏会受到局部环境影响。因此，行距测试既需要纯文本页，也需要包含典型元素的综合页面。下一章给出符号和公式，第三章给出表格、算法、代码与图形，可共同用于检查模板在真实科技论文场景中的稳定性。

% !TeX root = ../thesis.tex
\chapter{数学、物理符号与量和单位}
\label{chap:method}

\section{字母、数字与常用符号}
数学变量通常使用斜体，函数名、单位和说明性下标使用正体。小写希腊字母可显示为
$\alpha,\beta,\gamma,\delta,\epsilon,\varepsilon,\theta,\vartheta,\lambda,\mu,\nu,\xi,\pi,\rho,\sigma,\tau,\phi,\varphi,\chi,\psi,\omega$；大写形式包括
$\Gamma,\Delta,\Theta,\Lambda,\Xi,\Pi,\Sigma,\Phi,\Psi,\Omega$。常用集合与运算符包括
$\mathbb{R}$、$\mathbb{C}$、$\in$、$\notin$、$\subseteq$、$\cup$、$\cap$、$\forall$、$\exists$、$\nabla$、$\partial$、$\infty$、$\propto$、$\approx$、$\leq$ 和 $\geq$。

横线类字符外形相近，但编码和用途不同。键盘直接输入的 \verb|-| 是 U+002D，本义是连字符兼减号；中文破折号“——”由两个 U+2014 字符组成；英文范围常用的短破折号“--”由 LaTeX 排成 en dash；数学减号“$-$”则应使用 U+2212 对应的数学字形。U+2212 既用于二元减法运算，也用于表示负数的一元负号；其他运算符各有自己的字符或数学字形，并非所有数学符号都是 U+2212。

Word 正文中直接按键盘减号键，常会留下偏短、位置偏低的连字符，负号尤其容易被忽略。下面把同一表达式并列显示；错误行处于普通文本模式，两个短横线都是 U+002D，正确行处于数学模式，减法号和负号均按数学减号排版：
\begin{center}
  \begin{tabular}{cll}
    错误 & \textrm{a-b=-2} & 文本连字符 U+002D\\
    正确 & $a-b=-2$ & 数学减号 U+2212
  \end{tabular}
\end{center}
在 LaTeX 源码中写作 \verb|$a-b=-2$| 即可；虽然源码仍由键盘输入 ASCII 字符，数学模式会选择正确的减号字形和运算间距。

\section{行内公式和陈列公式}
行内公式应当作为句子的一部分，例如向量 $\boldsymbol{x}\in\mathbb{R}^{n}$、矩阵 $\boldsymbol{A}\in\mathbb{R}^{m\times n}$ 和转置 $\boldsymbol{x}^{\mathsf T}$。较长的推导应使用带编号的陈列公式：
\begin{equation}
  J(\boldsymbol{w})=\frac{1}{n}\sum_{i=1}^{n}\left(y_i-\boldsymbol{w}^{\mathsf T}\boldsymbol{x}_i\right)^2+\lambda\lVert\boldsymbol{w}\rVert_2^2.
  \label{eq:objective}
\end{equation}
式~\eqref{eq:objective} 同时展示求和、分式、范数、上下标、粗体变量及公式编号。分段函数、矩阵和微积分符号可写为
\begin{equation}
  f(x)=
  \begin{cases}
    \displaystyle \int_{0}^{x}\exp(-t^2)\,\mathrm{d}t, & x\geq 0,\\[4pt]
    \displaystyle -\int_{x}^{0}\exp(-t^2)\,\mathrm{d}t, & x<0,
  \end{cases}
  \qquad
  \boldsymbol{K}=\begin{bmatrix}k_{11}&k_{12}\\k_{21}&k_{22}\end{bmatrix}.
\end{equation}

多行公式使用对齐环境，使等号形成稳定的视觉轴：
\begin{align}
  \nabla\cdot\boldsymbol{q}+\frac{\partial \theta}{\partial t}&=s,\\
  \boldsymbol{q}&=-\boldsymbol{K}\nabla(h+z),\\
  \mathcal{L}\{f(t)\}(p)&=\int_{0}^{\infty} f(t)\exp(-pt)\,\mathrm{d}t.
\end{align}

学校撰写规范要求首次出现的物理量逐项注释时，破折号后的第一个字对齐。下面先给出一个完整公式，再用双长破折号逐项解释符号：
\begin{equation}
  \theta_{\mathrm f}=\frac{\mathrm d\varphi}{\mathrm d l}.
  \label{eq:torsion-angle}
\end{equation}
\par\noindent 式中，\par\nobreak
\begin{bjutformulaexplanation}
  \bjutexplain{$M_{\mathrm f}$}{试样断裂时的最大扭矩，单位为牛顿米（\si{\newton\meter}）；}
  \bjutexplain{$\theta_{\mathrm f}$}{试样断裂时单位长度上的相对扭转角，单位为弧度每毫米（\si{\radian\per\milli\meter}）；}
  \bjutexplain{$l$}{试样的有效长度，单位为毫米（\si{\milli\meter}）。}
\end{bjutformulaexplanation}
其中“式中，”与正文左边界对齐并单独占一行；变量注释另起一行，整体缩进两个汉字，双长破折号后的第一个字由表格列统一对齐。直接输入空格无法保证字体切换或换行后的稳定位置。若所在学科或期刊习惯把符号解释写成连续中文段落，也可采用“式中，$M_{\mathrm f}$ 为……；$\theta_{\mathrm f}$ 为……；$l$ 为……”的形式。两种样式均可使用，但同一篇论文宜保持一致。

\section{量、单位与乘号间距}
模板通过\mbox{siunitx}统一量值和单位格式，并把复合单位之间的乘号设为居中的“$\cdot$”。例如力矩可写为 \qty{18.6}{\newton\meter}，刚度可写为 \qty{18.6}{\newton\per\meter}，输出分别为“$\mathrm{N\cdot m}$”和“$\mathrm{N\cdot m^{-1}}$”。乘号左右的间距由宏包负责，不应手工插入空格。

其他典型量值包括温度 \qty{23.5}{\degreeCelsius}、压力 \qty{101.325}{\kilo\pascal}、速度 \qty{2.4e-3}{\meter\per\second}、密度 \qty{998}{\kilogram\per\cubic\meter}、角度 \ang{30} 以及范围 \qtyrange{10}{20}{\mega\pascal}。数字和单位之间保留不可断开的规范间距，单位符号使用正体，指数由单位结构自动生成。

\section{数学字体检查要点}
北京工业大学现行撰写规范把正文概括为“宋体，数字、字母等均用 Times New Roman”，但这一要求没有充分说明数学公式所需的专用字形及跨平台字体差异\cite{bjut2023guide}。Times New Roman 主要是西文正文字体，并不覆盖完整的数学运算符集合；复杂公式必然需要数学字体补足积分号、求和号、根号、伸缩括号等字形及其排版度量。因此，“宋体加 Times New Roman”可以描述正文的基本外观，却不能单独构成完整的科技写作字体方案。

Word 自带的 OfficeMath 通常以 Cambria Math 作为默认数学字体；MathType 7 的默认方案则组合使用 Times New Roman、Symbol 和 MT Extra，而不是只靠一种字体完成全部公式。macOS 自带字体清单中采用 Times、Times Roman 等名称，具体软件也可能随版本、安装来源和字体映射显示不同名称。因此，模板不把 Windows 中的字体菜单名称视为跨平台的唯一标准，而是分别控制西文正文与数学符号。当前实现以 Times New Roman 作为西文正文首选，以 \textsf{newtxmath} 提供 Times 风格的数学字母、运算符和度量；缺少专有字体时，再回退至兼容的 TeX Live 字体。判断输出是否合规时，应分别核对正文西文字体、数学字母和数学符号，不应仅凭某个字母的外观推断整份文档字体错误。

本模板的编排原则同时参考 GB/T 7713.1--2025 对学位论文组成和编排的规定，以及 GB/T 7714--2025 对参考文献著录和标引的规定\cite{gbt7713_1_2025,gbt7714_2025}。学校要求与国家标准发生差异时，提交版本仍应先确认研究生院和所在学院当年的明确要求。

\section{跨学科表达示例}
力学中的应力张量可写为 $\boldsymbol{\sigma}=\lambda\operatorname{tr}(\boldsymbol{\varepsilon})\boldsymbol{I}+2\mu\boldsymbol{\varepsilon}$；电路中的复阻抗可写为 $Z(\omega)=R+\mathrm{i}\omega L+1/(\mathrm{i}\omega C)$；统计量可写为 $\mathbb{E}[X]=\mu$ 和 $\operatorname{Var}(X)=\sigma^2$；化学反应则可用专用宏包表达为 \ce{CO2 + H2O <=> H2CO3}。这些例子可用于检查粗体张量、正体算子、虚数单位、黑板粗体和化学式中的上下标是否协调。

% !TeX root = ../thesis.tex
\chapter{图表、算法、代码与文献}
\label{chap:elements}

\section{三线表与数据表达}
表格应提供精确数据，正文负责解释趋势和意义。表~\ref{tab:comparison} 同时展示中文、数学符号、单位、有效数字和表下注释；表头不使用竖线，单位集中写在量名后，避免在每个数据单元中重复。

\begin{table}[htbp]
  \centering
  \caption{不同计算方案的结果比较}
  \label{tab:comparison}
  \begin{tabular}{lcccc}
    \toprule
    方案 & 网格数 $n$ & 时间步长 $\Delta t$/\si{\second} & 峰值压力/\si{\kilo\pascal} & 相对误差/\% \\
    \midrule
    基准方案 & 12,800 & 0.10 & 86.42 & --- \\
    加密网格 & 51,200 & 0.10 & 87.05 & 0.73 \\
    缩短步长 & 12,800 & 0.05 & 86.67 & 0.29 \\
    联合加密 & 51,200 & 0.05 & 87.12 & 0.81 \\
    \bottomrule
  \end{tabular}

  \vspace{2pt}
  \begin{minipage}{0.86\textwidth}\zihao{5}
  注：相对误差以基准方案为参照；“---”表示该项不适用，而不是缺失数据。
  \end{minipage}
\end{table}

\section{矢量图与双语题注}
流程图适合使用矢量格式，缩放时线条和文字不会像低分辨率位图一样模糊。图~\ref{fig:workflow} 由 TikZ 直接生成，展示输入、求解、验证和输出之间的关系。

\begin{figure}[htbp]
  \centering
  \begin{tikzpicture}[
    node distance=11mm,
    box/.style={draw,rounded corners=2pt,minimum width=30mm,minimum height=10mm,align=center,fill=black!3},
    arrow/.style={-{Latex[length=2.4mm]},thick}
  ]
    \node[box] (input) {数据与边界条件};
    \node[box,right=of input] (model) {模型与参数};
    \node[box,right=of model] (solve) {数值求解};
    \node[box,below=of model] (verify) {收敛性与误差检验};
    \node[box,right=of verify] (result) {结果解释与复现};
    \draw[arrow] (input) -- (model);
    \draw[arrow] (model) -- (solve);
    \draw[arrow] (solve) -- (result);
    \draw[arrow] (result) -- (verify);
    \draw[arrow] (verify) -- (model);
  \end{tikzpicture}
  \bicaption{可复现计算流程}{A reproducible computational workflow}
  \label{fig:workflow}
\end{figure}

\section{算法环境}
算法~\ref{alg:gradient} 展示输入、输出、循环、判断和数学表达式。伪代码用于解释方法逻辑；能够直接运行的实现应放在代码环境或独立源文件中。

\begin{algorithm}[htbp]
  \caption{带停止准则的梯度迭代}
  \label{alg:gradient}
  \begin{algorithmic}[1]
    \REQUIRE 初值 $\boldsymbol{w}_0$，步长 $\eta$，容差 $\varepsilon$
    \ENSURE 近似解 $\boldsymbol{w}$
    \STATE $k\leftarrow 0$
    \REPEAT
      \STATE $\boldsymbol{g}_k\leftarrow\nabla J(\boldsymbol{w}_k)$
      \STATE $\boldsymbol{w}_{k+1}\leftarrow\boldsymbol{w}_k-\eta\boldsymbol{g}_k$
      \STATE $k\leftarrow k+1$
    \UNTIL{$\lVert\boldsymbol{g}_{k-1}\rVert_2<\varepsilon$}
    \RETURN $\boldsymbol{w}_k$
  \end{algorithmic}
\end{algorithm}

\section{程序代码}
代码清单应使用等宽字体、保留缩进，并根据篇幅决定是否显示行号。清单~\ref{lst:python} 给出与算法~\ref{alg:gradient} 对应的 Python 实现。本地字体目录中存在更纱黑体等宽字体时，模板会优先使用该字体，从而同时支持中英文代码注释。

\begin{lstlisting}[language=Python,caption={梯度迭代的 Python 实现},label={lst:python}]
def gradient_descent(gradient, initial, step, tolerance, max_steps=1000):
    """Return an approximate stationary point."""
    value = initial
    for iteration in range(max_steps):
        direction = gradient(value)
        if norm(direction) < tolerance:
            return value, iteration
        value = value - step * direction
    raise RuntimeError("iteration did not converge")
\end{lstlisting}

\section{参考文献引用与著录类型}
LaTeX 的基本思想和文档组织方法可参见 Lamport 的专著~\cite{lamport1994latex}。连续编号引用可以同时包含英文期刊论文与中文规范性文件~\cite{rahardjo1995,zhang2024,bjut2023guide}。参考文献表还保留图书、国家标准、技术报告、政府网络文件等类型，用于观察作者数量、英文大小写、标准号、卷期页码、DOI 和网址的输出差异~\cite{fredlund1993,gbt7713_1_2025,gbt7714_2025,zgqhgb2025,stateCouncil2025}。

为使模板的参考文献样式得到充分展示，本示例列出数据库中的全部条目；正式论文不应使用 \verb|\nocite{*}|，而应只保留正文确实引用的文献。
\nocite{*}

% !TeX root = ../thesis.tex
\begin{conclusion}

\section*{主要结论}
本模板把封面、目录、正文、图表、公式、代码和参考文献纳入同一构建流程。正常编译完成后，PDF 和全部辅助文件位于 \texttt{build/}，源码目录保持整洁。

\section*{模板扩展}
模板相对于基础版本增加了可配置封面信息、本地字体优先与开源字体回退、Times 风格数学字体、复合单位居中点、双语题注、表格行距、代码字体以及适合 VS Code LaTeX Workshop 的构建规则。

\section*{后续核验}
学校格式要求可能更新，使用者应以提交当年的官方文件为准。建议在定稿前检查 PDF 页面尺寸、字体嵌入、目录页码、交叉引用、参考文献著录、图片分辨率和打印版行距；如果与 Word 基准页存在累计差异，应优先调整正文基线间距，而不是在各段落中零散插入空白。
\end{conclusion}

\makebib{bib/references}

\appendix
% \appendix 之后的章号自动切换为附录编号。
% !TeX root = ../thesis.tex
\chapter{编译与定稿检查清单}

附录适合收录较长推导、调查问卷、程序接口和补充数据。若论文不需要附录，可在主文件中注释对应的 \verb|\include| 行。定稿前建议依次完成以下检查：

\begin{enumerate}
  \item 使用 XeLaTeX 和 BibTeX 完成全量构建，确认日志中没有未定义引用；
  \item 用 \verb|pdffonts build/thesis.pdf| 检查要求的中英文字体是否嵌入；
  \item 以 100\% 比例比较 PDF 与学校 Word 样例中的页边距、正文行高和标题距离；
  \item 检查所有图表、公式和参考文献均在正文中被引用；
  \item 打印代表性页面，核对屏幕显示不易发现的线宽、灰度和字号差异。
\end{enumerate}

\backmatter
% !TeX root = ../thesis.tex
\begin{publication}
本部分只列入攻读学位期间已经公开发表、录用或按学校要求允许列出的成果，并保持作者顺序、成果状态和著录信息真实准确。没有相关成果时，应在 \texttt{thesis.tex} 中注释本部分，而不保留虚构条目。
\end{publication}

% !TeX root = ../thesis.tex
\begin{acknowledgement}
致谢用于说明对研究、写作、数据、设备或经费提供实质帮助的个人和机构。文字应简洁、真实，并在公开版本中避免披露未经同意的个人信息。资助项目应使用正式名称和编号，其位置以学校当年要求为准。
\end{acknowledgement}

\end{document}
